\documentclass[11pt,a4paper]{article}

\usepackage[utf8]{inputenc}
\usepackage[T1]{fontenc}
\usepackage[spanish]{babel}
\usepackage{lmodern}
\usepackage{amsmath,amssymb}
\usepackage{geometry}
\geometry{margin=2.4cm}
\usepackage{booktabs}
\usepackage{xcolor}
\usepackage{listings}
\usepackage{hyperref}
\usepackage{graphicx}
\usepackage{tikz}
\usetikzlibrary{shapes.geometric, arrows.meta, positioning, calc, fit, backgrounds}

\hypersetup{
    colorlinks=true,
    linkcolor=blue!60!black,
    urlcolor=blue!60!black,
    citecolor=blue!60!black
}

\definecolor{codebg}{rgb}{0.97,0.97,0.97}
\definecolor{codekw}{rgb}{0.10,0.10,0.60}
\definecolor{codestr}{rgb}{0.40,0.20,0.00}
\definecolor{codecom}{rgb}{0.30,0.50,0.30}

\lstdefinestyle{py}{
    backgroundcolor=\color{codebg},
    basicstyle=\ttfamily\footnotesize,
    keywordstyle=\color{codekw}\bfseries,
    stringstyle=\color{codestr},
    commentstyle=\color{codecom}\itshape,
    language=Python,
    breaklines=true,
    numbers=none,
    tabsize=4,
    showstringspaces=false,
    frame=single,
    framesep=4pt,
    xleftmargin=6pt,
    xrightmargin=6pt
}

\title{Extensión de Andromeda a \emph{gravitational lensing}\\
\large Métricas de lente, plano fuente y composición observacional}
\author{David Bamba\\\small Universidad Nacional de Colombia}
\date{\today}

\begin{document}
\maketitle

\begin{abstract}
Este documento registra la extensión del motor de \emph{ray tracing} de
Andromeda al régimen de \emph{gravitational lensing}. A partir del núcleo
numérico ya existente para agujeros negros (integrador Numba, paralelización
con \texttt{prange}, eventos, detector) se añaden (i) métricas de lente en
régimen de campo débil (SIS y SIE), (ii) una clase \texttt{SourcePlane} con
perfiles de luz 2D compilados (Gaussiano y Sérsic), (iii) un nuevo modo
\texttt{lensing} en el \emph{dispatcher} de trazado que intersecta los rayos
escapados con el plano fuente y (iv) una capa de composición visual
(galaxias de fondo, \emph{starfield}, ruido fotográfico, \emph{stretch}
\texttt{asinh}) para producir imágenes comparables a observaciones HST.
Se documentan las hipótesis físicas, el contrato de interfaz entre módulos,
los ejemplos implementados y las limitaciones conocidas.
\end{abstract}

\tableofcontents

\bigskip


\section{Motivación y alcance}

El núcleo de Andromeda fue diseñado para trazar geodésicas nulas en fondos
de Kerr y Schwarzschild (incluyendo variantes MOG y \emph{scalar-hair}) y
producir imágenes del disco de acreción. Toda la maquinaria numérica
---integrador adaptativo con detección de eventos, paralelización, detector
en el plano imagen--- es \textbf{agnóstica a la métrica}: cualquier
espaciotiempo que exponga los \emph{hooks} Numba
(\texttt{\_rhs\_nb}, \texttt{\_metric\_nb}, \texttt{\_omega\_nb}) puede
ser trazado sin tocar el \emph{hot path}.

La extensión a \emph{gravitational lensing} aprovecha esa neutralidad:
se añaden métricas de lente siguiendo el mismo contrato y se introduce un
nuevo tipo de objeto (\texttt{SourcePlane}) que reemplaza al disco de
acreción como destino de los rayos. Todo el código numérico existente
(integrador, eventos, paralelización) se reutiliza sin modificación.

\paragraph{Lo que NO se documenta aquí.}
Los detalles del integrador Numba, el \emph{dispatcher} de eventos, la
paralelización con \texttt{prange} y el contrato de los \emph{hooks}
métricos ya están cubiertos en \texttt{numba\_refactor.tex}. Aquí solo
se describe lo \emph{añadido} para lensing.


\section{Métricas de lente en campo débil}

\subsection{Forma general del \emph{line element}}

Las métricas implementadas parten de la forma isotrópica en campo débil
\begin{equation}
  ds^2 = -A(r,\theta)^2\, dt^2 + B(r,\theta)^2\,
         \bigl(dr^2 + r^2\,d\theta^2 + r^2\sin^2\theta\, d\phi^2\bigr),
  \label{eq:weak_field}
\end{equation}
con $A = 1 + \Phi$, $B = 1 - \Phi$ y $\Phi \ll 1$. Esto no es una solución
exacta de las ecuaciones de Einstein; es la forma canónica del lensing
gravitacional a primer orden en $\Phi$, que reproduce correctamente la
deflexión de la luz $\alpha = 2\Delta\Phi$ pero descarta términos
no-lineales en $\Phi$.

La unidad geometrizada es $G = c = 1$. No se introduce aún cosmología
FRW: las distancias $D_L$, $D_{LS}$ se fijan a mano en unidades de $M$.

\subsection{Singular Isothermal Sphere (SIS)}

Archivo: \texttt{scr/lens\_metrics/sis.py}.

Potencial esféricamente simétrico:
\begin{equation}
  \Phi(r) = 2\sigma_v^2 \ln\!\left(\frac{r}{r_{\rm ref}}\right),
\end{equation}
que reproduce el campo de una esfera isoterma singular con dispersión de
velocidades $\sigma_v$ (en unidades de $c$). La deflexión resultante es
constante en $b$ (parámetro de impacto):
\begin{equation}
  \alpha_{\rm SIS} = 4\pi\sigma_v^2,
\end{equation}
y produce un radio de Einstein $b_E = 4\pi\sigma_v^2\, D_{LS}$.

\subsection{Singular Isothermal Ellipsoid (SIE)}

Archivo: \texttt{scr/lens\_metrics/sie.py}.

Generalización ellipsoidal del SIS:
\begin{equation}
  \Phi(r,\theta) = 2\sigma_v^2 \ln\!\left(\frac{\xi}{r_{\rm ref}}\right),
  \qquad
  \xi = r\sqrt{\sin^2\theta + \frac{\cos^2\theta}{q_{\rm ax}^2}}.
\end{equation}

El parámetro $q_{\rm ax}\in(0,1]$ es el cociente de semiejes:
$q_{\rm ax} = 1$ recupera el SIS, $q_{\rm ax} < 1$ corresponde a un
elipsoide oblato respecto al eje $z$.

Las derivadas parciales (usadas por el \emph{hook} geodésico) son
\begin{align}
  \partial_r \Phi &= \frac{2\sigma_v^2}{r},\\[2pt]
  \partial_\theta \Phi &= 2\sigma_v^2\,
      \frac{\sin\theta\cos\theta\,(1 - 1/q_{\rm ax}^2)}
           {\sin^2\theta + \cos^2\theta/q_{\rm ax}^2}.
\end{align}

Se preserva la axisimetría en $\phi$ (el momento angular $k_\phi$ se
conserva), por lo que un rayo que parte en el plano ecuatorial permanece
allí.

\subsection{Contrato de interfaz}

Cada \texttt{LensMetric} expone exactamente los mismos atributos que
\texttt{BlackHole} (ver \texttt{numba\_refactor.tex}):
\begin{itemize}
  \item \texttt{\_rhs\_nb(q)} \ \ --- RHS geodésico compilado,
  \item \texttt{\_metric\_nb(x)} --- componentes de $g_{\mu\nu}$,
  \item \texttt{\_omega\_nb(\ldots)} --- velocidad angular nula
        (\emph{dummy} en lensing, requerido por el \emph{dispatcher}),
  \item \texttt{EH} --- pseudo-horizonte numérico
        (\texttt{r\_min}$\sim 10^{-2}$ actúa como corte del potencial
        logarítmico divergente).
\end{itemize}

Esto garantiza que \texttt{parallel.trace\_*} y el integrador no distinguen
entre un rayo en Kerr y un rayo en SIE: el mismo \texttt{solve\_ivp}
interno.


\section{Plano fuente y perfiles de luz}

\subsection{Clase \texttt{SourcePlane}}

Archivo: \texttt{scr/common/lens\_image.py}.

\begin{lstlisting}[style=py]
class SourcePlane:
    def __init__(self, D_LS, profile):
        self.D_LS    = float(D_LS)   # distancia lente-fuente
        self.profile = profile       # light profile 2D
        # duck-type para el dispatcher (no se usan en modo lensing)
        self._r_tbl = np.zeros(2, dtype=np.float64)
        self._I_tbl = np.zeros(2, dtype=np.float64)
        self.in_edge  = 0.0
        self.out_edge = 0.0
\end{lstlisting}

El plano se sitúa perpendicular al eje óptico en $x = -D_{LS}$. El
\emph{duck-type} con atributos heredados de \texttt{thin\_disk.structure}
permite que el \emph{dispatcher} de \texttt{parallel.py} acepte un
\texttt{SourcePlane} donde esperaba una estructura de disco, sin
ramificación adicional.

\subsection{Perfiles 2D compilados}

Archivo: \texttt{scr/sources/light\_profiles.py}.

Cada perfil empaca sus parámetros en un \texttt{float64[8]} y expone una
etiqueta entera \texttt{\_kind}. Un único dispatcher Numba evalúa el
brillo en $(x,y)$:

\begin{lstlisting}[style=py]
@njit(cache=True)
def _eval_profile_nb(xs, ys, kind, params):
    if kind == 0:   # Gaussian
        ...
    elif kind == 1: # Sersic
        ...
    return 0.0
\end{lstlisting}

\paragraph{Gaussiano.}
\begin{equation}
  I(x,y) = I_0\,\exp\!\left[-\frac{(x-x_0)^2 + (y-y_0)^2}{2\sigma^2}\right].
\end{equation}

\paragraph{Sérsic elíptico.}
\begin{equation}
  I(x,y) = I_e\,\exp\!\left[-b_n\!\left(\left(\frac{R}{R_e}\right)^{1/n}
  - 1\right)\right],
  \qquad
  R = \sqrt{x_r^2 + (y_r/q)^2},
\end{equation}
con $(x_r,y_r)$ rotado por el ángulo de posición \texttt{pa} alrededor
de $(x_0,y_0)$, $q = 1 - \varepsilon$ (axial ratio),
y $b_n$ aproximado por la expansión de Ciotti \& Bertin (1999)
\begin{equation}
  b_n \simeq 2n - \frac{1}{3} + \frac{4}{405\,n}
            + \frac{46}{25515\,n^2}.
\end{equation}

$n = 1$ corresponde a un disco exponencial; $n = 4$ a un perfil de
Vaucouleurs típico de elípticas.

\begin{figure}[!ht]
  \centering
  \includegraphics[width=\linewidth]{figs/profiles.pdf}
  \caption{Perfiles de luz 2D compilados
  (\texttt{scr/sources/light\_profiles.py}). Cada panel muestra el
  mapa de brillo $I(x,y)$ con \emph{stretch} asinh y, debajo, un corte
  unidimensional a $y = 0$ en escala log-$y$. La Gaussiana es la
  referencia circular; el Sérsic $n = 1$ reproduce un disco
  exponencial; el Sérsic $n = 4$ (de Vaucouleurs) concentra brillo
  hacia el centro con alas extendidas; la variante con
  $\varepsilon = 0{,}45$ es la fuente usada en \texttt{ex15} para
  producir arcos asimétricos. Generado por
  \texttt{docs/figs/gen\_profiles.py}.}
  \label{fig:profiles}
\end{figure}

\paragraph{Añadir un perfil nuevo.} Basta con (i) definir una nueva
etiqueta \texttt{KIND\_X}, (ii) añadir una rama en \texttt{\_eval\_profile\_nb}
y (iii) crear una clase con \texttt{\_kind} y \texttt{\_params} como
atributos. Ningún otro módulo necesita cambio.


\section{Modo \texttt{lensing} del \emph{dispatcher}}

El módulo \texttt{scr/common/parallel.py} acepta cuatro modos:
\begin{center}
\begin{tabular}{@{}ll@{}}
  \toprule
  \textbf{Código} & \textbf{Modo} \\
  \midrule
  0 & \texttt{no\_doppler} \\
  1 & \texttt{doppler} \\
  2 & \texttt{shadow} \\
  \textbf{3} & \textbf{\texttt{lensing}} (nuevo) \\
  \bottomrule
\end{tabular}
\end{center}

En modo \texttt{lensing}:

\begin{enumerate}
  \item El integrador traza el rayo \emph{backwards} desde el pixel del
        detector hasta que se dispara el evento \texttt{escape}
        ($r > r_{\rm escape}$).
  \item A partir del estado final
        $(x_{\rm esc}, k^\mu_{\rm esc})$ se extrapola en línea recta
        hasta el plano fuente: el rayo sale del campo gravitatorio y
        viaja en vacío. Se resuelve
        \begin{equation}
          x(\lambda) = -D_{LS},
        \end{equation}
        para $\lambda$ y se obtiene $(y_{\rm src}, z_{\rm src})$.
  \item Se evalúa el brillo
        $I = \text{\texttt{\_eval\_profile\_nb}}(y_{\rm src}, z_{\rm src},
        \text{\texttt{kind}}, \text{\texttt{params}})$
        y se asigna al pixel.
\end{enumerate}

\texttt{LensImage} (subclase de \texttt{Image}) es azúcar sintáctico para
fijar \texttt{mode="lensing"} al llamar \texttt{create\_image()}.

\begin{lstlisting}[style=py]
class LensImage(Image):
    def create_image(self, n_workers=None, chunksize=None):
        self._trace(mode="lensing",
                    n_workers=n_workers,
                    chunksize=chunksize)
\end{lstlisting}

\subsection{Geometría del detector}

El detector ya existente (\texttt{scr/detectors/image\_plane.py}) se usa
sin modificación. Para lensing conviene fijar $\iota = \pi/2$
(proyección paralela, observador en el infinito), de forma que cada pixel
sea un rayo con el mismo vector $\hat{k}$ pero diferente parámetro de
impacto $b = \sqrt{\alpha^2 + \beta^2}$.

\begin{figure}[!ht]
  \centering
  \includegraphics[width=\linewidth]{figs/geometry.pdf}
  \caption{Geometría del pipeline de lensing. El plano fuente (izquierda,
  $x = -D_{LS}$) contiene la distribución de brillo $I(y,z)$; la lente
  (centro, naranja) curva las geodésicas nulas; el detector (derecha,
  $x = +D_L$) registra un pixel $(\alpha,\beta)$ por cada rayo trazado
  \emph{hacia atrás}. Las trayectorias azules y rojas corresponden a
  geodésicas nulas reales (SIS, $\sigma_v = 0{,}03\,c$); las líneas
  grises discontinuas son el contrafactual sin lensing; las
  dos imágenes a $\pm b_E$ forman el anillo de Einstein. Generado por
  \texttt{docs/figs/gen\_geometry.py}.}
  \label{fig:geometry}
\end{figure}


\section{Composición de imagen observacional}

El contenido anterior basta para producir un \emph{arco bruto} del rayo
lensado. Para obtener imágenes con aspecto de observación real se añadió
una capa de post-proceso en \texttt{scr/common/visual.py}:

\begin{description}
  \item[\texttt{project\_profile\_on\_detector}] evalúa directamente un
        perfil 2D en el plano imagen, sin \emph{ray tracing}. Se usa para
        la luz propia de la galaxia-lente (que no se lensa a sí misma) y
        para galaxias de fondo proyectadas.
  \item[\texttt{add\_starfield}] genera $N$ estrellas puntuales con PSF
        Gaussiana, flujos log-uniformes y posiciones aleatorias.
  \item[\texttt{add\_photographic\_noise}] aplica ruido Poisson
        (escalado por \texttt{poisson\_scale}) más ruido Gaussiano de
        lectura. No modela \emph{cosmic rays} ni \emph{dark current}.
  \item[\texttt{compose\_rgb}] combina capas monocromáticas en RGB
        aplicando colores prefijados por componente.
  \item[\texttt{asinh\_stretch}] \emph{stretch} percepción-lineal al
        estilo HST:
        \begin{equation}
          I_{\rm out} = \frac{\mathrm{asinh}(I/s)}{\mathrm{asinh}(I_{\max}/s)},
        \end{equation}
        con \emph{softening} $s$ y máximo tomado del percentil 99.7 de
        la luminancia.
\end{description}


\section{Pipeline completo}

\begin{tikzpicture}[
  node distance=6.5mm,
  every node/.style={rectangle, rounded corners=2pt, draw=black!65,
                     inner sep=4pt, font=\footnotesize, align=center,
                     minimum width=3.6cm},
  arr/.style={-{Stealth[length=1.6mm]}, thick, black!70}
]
  \node (metric) [fill=blue!6]  {\texttt{LensMetric} (SIS / SIE)\\
                                  \scriptsize hooks Numba};
  \node (det)    [below=of metric, fill=blue!6]
                                 {\texttt{detector}\\\scriptsize $\iota, D_L$, pixeles};
  \node (phot)   [below=of det, fill=blue!6]
                                 {\texttt{create\_photons()}\\\scriptsize $k^\mu_0$ por pixel};
  \node (trace)  [below=of phot, fill=orange!12]
                                 {\texttt{parallel.trace\_*}\\\scriptsize mode=\texttt{lensing}};
  \node (src)    [below=of trace, fill=green!8]
                                 {\texttt{SourcePlane} + \texttt{light\_profiles}\\
                                  \scriptsize Sérsic / Gaussian};
  \node (vis)    [below=of src, fill=red!8]
                                 {\texttt{visual.*}\\
                                  \scriptsize starfield, noise, asinh};
  \node (out)    [below=of vis]  {PNG + NPY};

  \draw[arr] (metric) -- (det);
  \draw[arr] (det)    -- (phot);
  \draw[arr] (phot)   -- (trace);
  \draw[arr] (trace)  -- (src);
  \draw[arr] (src)    -- (vis);
  \draw[arr] (vis)    -- (out);
\end{tikzpicture}


\section{Ejemplos implementados}

Todos los ejemplos viven en la raíz del repositorio y guardan sus salidas
en \texttt{images/lensing/} y \texttt{images\_data/lensing/}.

\begin{center}
\renewcommand{\arraystretch}{1.25}
\begin{tabular}{@{}p{3.4cm}p{12.5cm}@{}}
  \toprule
  \textbf{Ejemplo} & \textbf{Qué valida / muestra} \\
  \midrule
  \texttt{ex9.einstein\_ring} &
  Schwarzschild puntual + fuente Gaussiana en el eje $\Rightarrow$ anillo
  de Einstein. Se verifica
  $b_{\rm ring} = 2\sqrt{M\,D_{LS}}$. \\
  \texttt{ex10.cosmic\_horseshoe} &
  SIS + Sérsic desplazado + composición RGB con \emph{starfield} suave.
  Produce un arco tipo ``Cosmic Horseshoe'' (LRG 3-757). \\
  \texttt{ex11.lensing\_trajectory} &
  Trazado directo de geodésicas nulas en SIS; diagrama con ángulos de
  deflexión y parámetros de impacto. \\
  \texttt{ex12.kerr\_lensing} &
  Kerr ($a = 0.9$) vs Schwarzschild lado a lado; visualiza la asimetría
  por \emph{frame-dragging} en régimen de campo fuerte. \\
  \texttt{ex13.einstein\_cross} &
  SIS vs SIE ($q = 0.55$): anillo simétrico contra configuración de
  imágenes múltiples (cruz de Einstein). \\
  \texttt{ex14.realistic\_lensing} &
  Primera capa de realismo observacional: múltiples fuentes lensadas,
  galaxias de fondo no-lensadas, \emph{starfield}, ruido fotográfico. \\
  \texttt{ex15.broken\_ring\_lensing} &
  SIE + fuente primaria fuertemente elíptica ($\varepsilon = 0{,}45$);
  arcos asimétricos con ``puentes'' entre imágenes múltiples. \\
  \bottomrule
\end{tabular}
\end{center}


\section{Hipótesis físicas y validez}

\begin{figure}[!ht]
  \centering
  \includegraphics[width=\linewidth]{figs/trajectory_comparison.pdf}
  \caption{Mismo integrador, dos regímenes. Izquierda: Schwarzschild
  ($M = 1$), régimen de campo fuerte; los rayos con parámetro de impacto
  cerca de $b_{\rm crit} = 3\sqrt{3}\,M$ rodean la esfera de fotones y
  algunos terminan cayendo al horizonte. Derecha: lente SIS
  galáctico ($\sigma_v = 0{,}03\,c$); los rayos se deflectan suavemente
  en un régimen de campo débil $|\Phi| \ll 1$ y todos escapan. Ambos
  casos usan el mismo integrador (RK45/DOP853) y el mismo sistema de
  eventos, cambiando únicamente el objeto métrica. Generado por
  \texttt{docs/figs/gen\_trajectory\_comparison.py}.}
  \label{fig:trajcomp}
\end{figure}

\begin{itemize}
  \item \textbf{Campo débil.} La métrica (\ref{eq:weak_field}) es
        estrictamente válida para $|\Phi| \ll 1$. En los ejemplos se usa
        $\sigma_v \sim 0{,}028\,c$ (velocidades infladas respecto a
        galaxias reales, $\sim 200$\,km/s $\approx 10^{-3}\,c$) solo
        para obtener radios de Einstein visualmente representativos en
        unidades geometrizadas. La deflexión angular reproducida sigue
        siendo correcta a primer orden, pero el régimen no es
        estrictamente ``weak'' cuando se fuerzan estos valores.
  \item \textbf{Axisimetría.} Todas las métricas de lente conservan
        $k_\phi$. Lentes no-axisimétricas generales (p.ej.\ con
        \emph{external shear} en dirección arbitraria) requieren romper
        esta simetría y no están implementadas.
  \item \textbf{Un único plano fuente.} No hay \emph{multi-plane
        lensing}: la intersección se hace con un solo
        \texttt{SourcePlane} a distancia $D_{LS}$ fija.
  \item \textbf{Plano fuente plano.} La fuente no se extiende en
        \emph{redshift}: todo el perfil vive en $x = -D_{LS}$ exactamente.
  \item \textbf{Cosmología ausente.} No se han introducido distancias
        angulares de diámetro FRW; $D_L$ y $D_{LS}$ son parámetros libres
        en unidades de $M$. Esto impide convertir a arcsec o kpc.
  \item \textbf{Galaxias de fondo no lensadas.} En los composites
        (\texttt{ex10}, \texttt{ex14}, \texttt{ex15}) las galaxias de
        fondo se proyectan directamente; en la realidad también sufren
        convergencia y \emph{shear}, aunque en régimen débil.
  \item \textbf{Óptica del detector.} No hay PSF de Hubble ni patrón de
        difracción. La resolución está limitada solo por el muestreo
        angular del plano imagen.
  \item \textbf{Ruido.} Modelo simple Poisson + Gaussiano. No se
        modelan \emph{dark current}, \emph{cosmic rays}, \emph{flat-field
        residuals} ni \emph{charge transfer inefficiency}.
  \item \textbf{Colores.} Arbitrarios, no derivados de SED por
        \emph{redshift} por filtro; son una convención visual.
\end{itemize}


\section{Qué está validado y qué falta validar}

\paragraph{Validado numéricamente.}
\begin{itemize}
  \item Radio de Einstein de Schwarzschild
        $b_E = 2\sqrt{M\,D_{LS}}$ (\texttt{ex9}).
  \item Deflexión constante en SIS
        $\alpha = 4\pi\sigma_v^2$ (verificación en \texttt{ex11}).
  \item Hamilton-constraint en cada paso del integrador
        (\texttt{Image.verify\_Hamiltonian()}, heredado del trazado BH).
\end{itemize}

\begin{figure}[!ht]
  \centering
  \includegraphics[width=\linewidth]{figs/deflection.pdf}
  \caption{Validación cuantitativa de la deflexión angular. Izquierda:
  en Schwarzschild la deflexión $\alpha(b)$ medida por el ray-tracer
  (puntos rojos) sigue la predicción analítica a primer orden
  $\alpha = 4M/b$ (línea negra); se evitan parámetros de impacto cerca
  de $b_{\rm crit} = 3\sqrt{3}\,M$ donde los términos no-lineales
  dominan. Derecha: en SIS la deflexión es independiente del parámetro
  de impacto ($\alpha = 4\pi\sigma_v^2$, línea negra) y los puntos
  numéricos (azules) reproducen el plateau. El error relativo mediano
  se muestra en cada panel. Generado por
  \texttt{docs/figs/gen\_deflection.py}.}
  \label{fig:deflection}
\end{figure}

\paragraph{Pendiente de validación (propuesta).}
\begin{itemize}
  \item \textbf{Test de topología de imágenes múltiples}
        (\emph{caustic validation}): con SIE fijo, variar solo el
        \emph{offset} de la fuente para reproducir las transiciones
        $4 \to 3 \to 2$ imágenes al cruzar las cáusticas tangencial y
        radial. Si el kernel numérico reproduce las cuatro topologías
        con un único parámetro, el trazado en lensing está validado más
        allá de la simple posición del anillo.
  \item \textbf{Magnificación}: recuperar $\mu = 1/\det(A)$ por
        diferenciación numérica del mapa $\beta(\theta)$ y comparar con
        fórmulas analíticas para SIS/SIE (no implementado aún).
\end{itemize}


\section{Resumen de archivos añadidos / modificados}

\begin{center}
\renewcommand{\arraystretch}{1.2}
\begin{tabular}{@{}p{5.4cm}p{10.5cm}@{}}
  \toprule
  \textbf{Archivo} & \textbf{Rol} \\
  \midrule
  \texttt{scr/lens\_metrics/sis.py} & Métrica SIS (nuevo). \\
  \texttt{scr/lens\_metrics/sie.py} & Métrica SIE (nuevo). \\
  \texttt{scr/sources/light\_profiles.py} &
    Gaussian + Sérsic + dispatcher Numba (nuevo). \\
  \texttt{scr/common/lens\_image.py} &
    \texttt{SourcePlane} y \texttt{LensImage} (nuevo). \\
  \texttt{scr/common/visual.py} &
    \emph{starfield}, ruido, RGB, \emph{asinh stretch} (nuevo). \\
  \texttt{scr/common/parallel.py} &
    Añadido modo \texttt{lensing} al \emph{dispatcher} (modificado). \\
  \texttt{ex9 -- ex15} &
    Ejemplos de validación y composición observacional (nuevo). \\
  \bottomrule
\end{tabular}
\end{center}

Ninguno de los archivos del núcleo numérico BH
(\texttt{integrator.py}, \texttt{\_solvers.py}, \texttt{\_numba\_kernels.py})
fue modificado: toda la extensión es aditiva.


\section{Siguientes pasos}

La hoja de ruta para aproximar aún más las imágenes a observaciones
reales está descrita en \texttt{LENSING.md} y se resume en:
\begin{enumerate}
  \item Introducir cosmología FRW ($D_L$, $D_S$, $D_{LS}$ por
        \emph{redshift}; conversión a arcsec / kpc).
  \item Añadir métricas adicionales: NFW, PEMD/EPL, \emph{external
        convergence} y \emph{shear}.
  \item Calcular curvas críticas, cáusticas, mapas de magnificación y
        \emph{time-delay}.
  \item Convolución con PSF real (HST F814W/F160W) y modelo de ruido
        observacional completo.
  \item Fuentes extendidas no paramétricas (imágenes reales como perfil
        de fuente).
\end{enumerate}


\end{document}
